\section{Alignments and a displacement inequality}\label{sec:displacement}

The word construction will be measured first by insertion--deletion
distance. Its advantage is that every possible comparison has a concrete
description as an increasing matching. We then state the analytic
inequality that will be applied at each level of the construction.

\subsection{Distances from increasing matchings}

Write $\ID(x,y)$ for the minimum number of single-symbol insertions and
deletions transforming $x$ into $y$. An increasing matching is a list of
pairs of equal-symbol positions, strictly increasing in both coordinates.
Let $\LCS(x,y)$ be its maximum size. These are the classical trace and
subsequence descriptions of string correction
\cite[Sections~3 and~5]{WagnerFischer1974}.

\begin{lemma}\label{lem:alignments}
For arbitrary finite strings, including the empty string,
\begin{equation}\label{eq:metric-comparison}
 \ED(x,y)\le\ID(x,y)\le2\ED(x,y),
 \qquad
 \ID(x,y)=|x|+|y|-2\LCS(x,y).
\end{equation}
For two strings of the same length $n$, define their one-sided deficit by
\[
 r_n(x,y)=n-\LCS(x,y)=\tfrac12\ID(x,y).
\]
A matching of size $n-E$ leaves exactly $E$ unmatched positions on each
side.
\end{lemma}

\begin{proof}
An insertion--deletion script is an edit script, and replacing each
substitution by one deletion and one insertion proves the two
inequalities. Deleting the unmatched source positions in an increasing
matching and inserting the unmatched target positions realizes its
total loss $|x|+|y|-2|M|$. Conversely, the original positions that
survive an insertion--deletion script form an increasing matching; its
loss is at most the number of operations. Minimizing the loss gives the
identity. The equal-length assertions follow directly.
\end{proof}

\subsection{The comparison to be proved}

For a finite abelian group $G$ and a map $H:G\to\ell_1$, set
\begin{equation}\label{eq:variation}
 V_H(g)=\E_{z\in G}\norm{H(z+g)-H(z)}.
\end{equation}
Every expectation over a finite set uses the uniform probability measure.
The following proposition compares a shift in one component with a
simultaneous shift and with shifts averaged inside the separate
components. In the word construction, the simultaneous shift will move
the row list by one; the component averages will be bounded at the
preceding level.

\begin{proposition}[Prime-product displacement inequality]
\label{prop:displacement}\label{fourier:displacement}
Let $b,h\ge1$ be integers, let $P\ge1$, and let $G_1,\ldots,G_b$ be finite
abelian groups. Suppose $v_j\in G_j$ has order $p_j$, where
$p_1,\ldots,p_b$ are distinct primes in $[P,2P]$. In
$G=\prod_{j=1}^bG_j$, use the same notation $v_j$ for the element acting
only in coordinate $j$, and put $\tau=\sum_{j=1}^b v_j$. Then every
$H:G\to\ell_1$ and every integer $u$ satisfy
\begin{equation}\label{eq:prime-product-bound}
 V_H(uv_1)
 \le (2P)^{2h}V_H(\tau)
 +\frac2h\sum_{j=1}^b\E_{a\in\Z/p_j\Z}V_H(av_j).
\end{equation}
\end{proposition}

We prove the proposition after recording two elementary forms of the
cut and Fourier methods. The cut representation is standard
\cite[Section~4]{LinialLondonRabinovich1995}. Its summability statement
below permits an unrestricted sequence-space target.

\begin{lemma}[Finite cut decomposition]\label{lem:cuts}
Let $X$ be finite and $H:X\to\ell_1$. There are finite numbers $c_A\ge0$,
indexed by the nonempty proper subsets of $X$, such that
\begin{equation}\label{eq:cuts}
 \norm{H(x)-H(y)}
 =\sum_{\varnothing\ne A\subsetneq X}
 c_A|\mathbf1_A(x)-\mathbf1_A(y)|
 \qquad(x,y\in X).
\end{equation}
Moreover,
\[
 \sum_{\varnothing\ne A\subsetneq X}c_A
 \le\sum_{\{x,y\}\subset X,\ x\ne y}\norm{H(x)-H(y)},
\]
where the right sum is over unordered pairs.
\end{lemma}

\begin{proof}
For $|X|\le1$ all the sums are empty. Otherwise, for coordinate $n$ of
$H$, list its distinct values on $X$ as
$t_{n,1}<\cdots<t_{n,k_n}$. Put
\[
 A_{n,r}=\{x:H_n(x)>t_{n,r}\},\qquad
 a_{n,r}=t_{n,r+1}-t_{n,r}\quad(1\le r<k_n).
\]
The consecutive gaps telescope:
\[
 |H_n(x)-H_n(y)|
 =\sum_{r=1}^{k_n-1}a_{n,r}
 |\mathbf1_{A_{n,r}}(x)-\mathbf1_{A_{n,r}}(y)|.
\]
Group all terms with the same cut by setting
$c_A=\sum_{n,r:\,A_{n,r}=A}a_{n,r}$, initially allowing the value
$+\infty$. Nonnegativity allows the coordinate series to be interchanged
with the finite sum over unordered pairs. Hence
\[
 \sum_{\varnothing\ne A\subsetneq X}c_A|A|\,|X\setminus A|
 =\sum_{\{x,y\}\subset X,\ x\ne y}\sum_{n=1}^{\infty}
 |H_n(x)-H_n(y)|
 =\sum_{\{x,y\}\subset X,\ x\ne y}\norm{H(x)-H(y)}<\infty.
\]
Each indexed cut separates at least one pair, so
$|A|\,|X\setminus A|\ge1$. The displayed equality proves that all the
weights are finite and bounds their total. Summing the coordinate
identity and grouping its nonnegative terms now gives \eqref{eq:cuts}.
\end{proof}

\begin{lemma}[Finite Fourier displacement identity]\label{lem:fourier}
Let $T=\prod_{j=1}^b\Z/p_j\Z$, where the integers $p_j$ are at least two.
For $\lambda\in T$, set
\[
 \chi_\lambda(t)=
 \exp\!\left(2\pi\mathrm i\sum_{j=1}^b\frac{\lambda_jt_j}{p_j}\right),
 \qquad
 \widehat F(\lambda)=\E_{t\in T}F(t)\overline{\chi_\lambda(t)}.
\]
For every $F:T\to\mathbb C$ and $l\in T$,
\begin{equation}\label{eq:fourier-displacement}
 \E_{t\in T}|F(t+l)-F(t)|^2
 =\sum_{\lambda\in T}|\widehat F(\lambda)|^2
 4\sin^2\!\left(\pi\sum_{j=1}^b\frac{l_j\lambda_j}{p_j}\right).
\end{equation}
The squared sine is independent of the chosen integer representatives.
\end{lemma}

\begin{proof}
The geometric-series identity
\[
 \frac1p\sum_{a=0}^{p-1}e^{2\pi\mathrm iqa/p}
 =\begin{cases}
 1,&p\mid q,\\
 0,&p\nmid q
 \end{cases}
\]
shows, coordinate by coordinate, that the $|T|$ characters are
orthonormal for the inner product
$\langle F_1,F_2\rangle=\E_tF_1(t)\overline{F_2(t)}$.
They form a basis because the function space has dimension $|T|$.
Thus $F(t)=\sum_\lambda\widehat F(\lambda)\chi_\lambda(t)$, and the
squared norm equals the sum of squared Fourier coefficients. Apply this
to
\[
 F(t+l)-F(t)
 =\sum_\lambda\widehat F(\lambda)
       (\chi_\lambda(l)-1)\chi_\lambda(t)
\]
and use $|e^{2\pi\mathrm i\theta}-1|^2=4\sin^2(\pi\theta)$.
\end{proof}

\begin{proof}[Proof of Proposition~\ref{prop:displacement}]
By Lemma~\ref{lem:cuts}, each displacement of $H$ is a nonnegative
linear combination of displacements of Boolean cut maps. It is enough
to prove \eqref{eq:prime-product-bound} for one such map
$H=\mathbf1_A$.

Let $K=\prod_j\langle v_j\rangle\subseteq G$. Every shift in the desired
inequality preserves every coset of $K$. On a fixed coset, choose a
representative $c$ and identify
$T=\prod_j\Z/p_j\Z$ with that coset through
$t\mapsto c+\sum_jt_jv_j$. The cut becomes
$F(t)=\mathbf1_A(c+\sum_jt_jv_j)$. For Boolean values, absolute
difference equals squared difference. With
$D_F(l)=\E_t|F(t+l)-F(t)|^2$, it remains to prove
\begin{equation}\label{eq:coset-bound}
 D_F(ue_1)\le(2P)^{2h}D_F(\mathbf1)
 +\frac2h\sum_{j=1}^b\E_{a\in\Z/p_j\Z}D_F(ae_j),
\end{equation}
where $\mathbf1=(1,\ldots,1)$ and $e_j$ is the $j$th coordinate vector.

We compare the multipliers in Lemma~\ref{lem:fourier} one frequency at
a time. Fix $\lambda\in T$, write
$S=\{j:\lambda_j\ne0\}$ and $r=|S|$. When $\lambda_1=0$, the
left multiplier vanishes. Assume therefore that $\lambda_1\ne0$.
The left multiplier is at most four.

First suppose $r<h$. Choose $0\le\lambda_j<p_j$ and put
\[
 q=\prod_{j\in S}p_j,\qquad
 \theta=\sum_{j\in S}\frac{\lambda_j}{p_j}.
\]
For any $i\in S$, distinctness of the primes gives
$q\theta\equiv\lambda_i(q/p_i)\not\equiv0\pmod{p_i}$.
Thus $\theta$ is not an integer. Since its denominator divides $q$,
its distance $\delta$ from the nearest integer satisfies
\[
 \frac12\ge\delta\ge\frac1q
 \ge(2P)^{-r}\ge(2P)^{-h}.
\]
Concavity of sine on $[0,\pi/2]$ gives
$\sin(\pi\delta)\ge2\delta$. Consequently the simultaneous-shift
multiplier on the right of \eqref{eq:coset-bound} is at least
\[
 (2P)^{2h}4\sin^2(\pi\theta)
 =(2P)^{2h}4\sin^2(\pi\delta)\ge16,
\]
which bounds the left multiplier.

Now suppose $r\ge h$. For each coordinate, the same geometric-series
identity gives
\[
 \E_{a\in\Z/p_j\Z}4\sin^2(\pi a\lambda_j/p_j)
 =2-2\operatorname{Re}\E_a e^{2\pi\mathrm ia\lambda_j/p_j}
 =\begin{cases}2,&j\in S,\\0,&j\notin S.\end{cases}
\]
The component terms on the right therefore have combined multiplier
$(2/h)2r\ge4$, again bounding the left. This case includes $h=1$.
Multiplication by $|\widehat F(\lambda)|^2$ and summation over
$\lambda$ proves \eqref{eq:coset-bound}. Averaging over the equal-size
cosets gives the uniform average over $G$. Finally, summing the
nonnegative cut weights proves the proposition for $H$.
\end{proof}

The two cases explain the role of distinct periods. A low-support
frequency pays for a nonintegral simultaneous phase; a high-support
frequency pays for many component variations. The next section turns
these two payments into an iterated obstruction.
